\chapter{Faget som ikkje kan seie nei}

\begin{motto}
Palantir er eit tenestedesignprosjekt.\\ Faget har ikkje språket til å seie noko anna.
\end{motto}

Alt eg har sagt så langt kunne ein lese som ein akademisk krangel — interessant for dei i faget, utan følgjer for verda. Dette kapitlet handlar om følgjene. For mangelen på fundament er ikkje berre ein epistemisk pinleg detalj. Han har ein moralsk konsekvens, og konsekvensen er at faget strukturelt ikkje kan seie nei til noko. Ikkje fordi designarar er dårlege menneske; dei er som folk flest. Men fordi eit fag utan ein presis modell av kva ein form gjer — kva agentar ho rører, kva trykk ho ber, kva ho legg til side — manglar den interne målestokken som ville late det skilje forma som ber verda med seg frå forma som sviktar langs alt han ikkje såg. Og det som ikkje kan målast, kan ikkje avvisast.

\section{Den eindimensjonale forma}

Tenk på den enklaste designfeilen, og forstørr han. Éin eingongskopp tek tolv sekund å lage, tjue minutt å bruke, og fire hundre år å bryte ned. Det er ikkje dårleg design. Det er presis design — presis for éin akse. Koppen løyser nett det problemet han vart beden om å løyse, og sviktar langs alle dei andre: nedbrytinga, vatnet, jorda, dyra, generasjonane som arvar avfallet. Svikten er ikkje ein biverknad. Ho er signaturen til ein formgjevingsprosess som representerte éin dimensjon og var blind for resten.

Skaler opp, og du finn det same mønsteret overalt der noko er formgjeve utan eit fundament som tvingar fram fleire aksar. Energisystemet optimert for billeg kraft, blindt for atmosfæren. Matsystemet optimert for kalori per hektar, blindt for jordsmonnet. Plattformene optimerte for merksemd, blinde for borna sine sinn. Éin akse dominerer; dei andre forsvinn; forma sviktar langs det som forsvann; og svikten ber signaturen av blindskapen som laga han, like sikkert som eit fossil ber forma til dyret. Faget har ingen reiskap for å sjå denne signaturen på førehand, fordi det ikkje har eit omgrep om forma som ein posisjon i eit fleirdimensjonalt rom der nokre aksar er representerte og andre lagde til side. Det har estetisk intuisjon, stilhistorie og handverksminne. Det har aldri vore nok. No, med verktøya vi har, er det farleg.

\section{Når objektet er eit menneske}

Tenestedesign lagar og optimaliserer system for målretting, profilering og deportasjon, og kallar det funksjonelt. Palantir er eit tenestedesignprosjekt. Kunden er staten; «brukaren», i den grad ordet tyder noko, er den som vert overvaka, skåra, sortert, fjerna. Systemet vart bygt for etterretning, selt til politi og grensekontroll, eksportert til militære operasjonar verda over. Det vert brukt til systematisk terrorisering av innvandrarfamiliar; born ned i null år vert registrerte som mål i ein database\kref{20}. Variantar av same arkitekturen styrer målval i operasjonar som oppfyller folkerettslege definisjonar av massedrap; kvart mål er eit punkt i eit formrom optimert for éin akse\kref{21}. Systema er funksjonelle. Dei er presist funksjonelle langs den eine aksen dei vart bedne om å optimalisere.

Her ser du kvifor det manglande fundamentet er eit moralsk og ikkje berre eit teoretisk problem. Det gjeldande rammeverket i faget har inga grense for korleis ein kan handsame eit objekt, uavhengig av kva objektet er. Så lenge handsaminga kan kallast funksjonell, er ho innanfor — òg når objektet er levande, òg når objektet er eit menneske. Eit fag som ikkje kan seie kva ein agent er, kva eit seleksjonstrykk er, kva som vert lagt til side, har ingen stad å stå når det skal seie at \emph{dette} skal ikkje lagast. Det manglar setninga. Og det som manglar setninga, endar med å bygge tingen.

\section{Etikken som vart bolta på}

Faget sitt svar på dette er å leggje til eit emne. Det finst no etikk på pensum, berekraft på pensum, klima på pensum; det er pålagt, og det er meint vel. Men sjå kva det inneber at etikken er eit \emph{tillegg}. Det inneber at faget sjølv — modellen av kva design er — er etisk nøytralt, og at moralen kjem utanfrå og vert lagd oppå. Det er nett feil veg. I kvart fag som faktisk har eit fundament foreinleg med korleis verda heng saman, fell etikken ut av faget sjølv. Økologien treng ikkje eit påbolta etikkemne for å vite at eit inngrep i éin node forplantar seg gjennom heile nettet; det \emph{er} økologien. Medisinen treng ikkje eit eige emne for å vite at kroppen ikkje kan skjerast i delar utan tap; det er fysiologien. Forskingsetikken i desse faga følgjer direkte av kunnskapen deira: du får ikkje gjere kva du vil, fordi handlinga di ikkje stoppar der du sluttar å sjå.

Eit kvart fundament foreinleg med samtidas forståing av økologi, biologi og kognisjon ville uunngåeleg produsere ein etikk som la avgrensingar på designaren — ikkje som eit tillegg, men som ein konsekvens. At designfaget i staden må bolte etikken på utanfrå, er sjølve beviset på at faget sin modell av seg sjølv ikkje er foreinleg med korleis verda heng saman. Faget har valt å ikkje vite det nabofaga veit, ikkje fordi kunnskapen er utilgjengeleg — han ligg hjå kvart nabofag — men fordi eit fag som tek den kunnskapen innover seg, ville måtte slutte å gjere mykje av det faget gjer i dag. Det manglande fundamentet er ikkje ein uhell. Det er ein bekvem tilstand. Det held faget fritt frå å måtte seie nei.

\section{Skulda er kollektiv, og difor strukturell}

Lat meg vere tydeleg om kven dette råkar og kven det ikkje råkar. Eg skuldar ikkje den einskilde designaren som tok eit oppdrag for å leve. Eit individuelt nei tapar berre oppdraget; det neste byrået tek det. Skulda er ikkje individuell. Ho er kollektiv og strukturell, og difor kan ho berre kurerast kollektivt og strukturelt. Eit fag som hadde eit fundament, kunne gjort det individuelle nei om til eit fagleg nei — slik medisinen kan strippe ein lege for retten til å praktisere, slik ingeniørstanden kan utelukke for grov aktløyse. Designfaget har ingen slik mekanisme, fordi ein slik mekanisme krev ein standard å utelukke for brot på, og faget har ingen standard. Det har ingen ting å bli utelukka frå. Du kan ikkje miste autorisasjonen i eit fag som ikkje har ein.

Og difor er dette kapitlet det tyngste leddet i klagemålet. Dei føregåande kapitla viste at faget ikkje kan grunngje seg sjølv. Dette viser kva det manglande fundamentet kostar: ein profesjon som byggjer det verste like villig som det beste, fordi han manglar språket til å skilje dei. Faget likar å sjå på seg sjølv som ei kraft for det gode i verda — menneskesentrert, empatisk, problemløysande. Men ein profesjon utan evne til å seie nei er ikkje ei kraft for det gode. Han er ei kraft, ubunden, til disposisjon for den som betalar, og resultatet ber signaturen til kven som helst som leigde han.

\vignett

Korleis kan eit fag i denne tilstanden halde fram, tiår etter tiår, utan at samanbrotet vert openbert for alle? Svaret ligg i institusjonen — i akkrediteringa, prisane, konferansane, heile maskineriet som produserer ein illusjon av at her finst eit fag med standardar. Det er sjølvbedraget, og det er neste kapittel.
